% =====================================================================
\section{The Paper's Model: Blocks 1--6}
\label{sec:paper}
% =====================================================================

This section follows \cref{fig:flowpaper} from the plant to the adaptive law. All numbers in the toy examples use the
paper's Part~1 initial state at $t=0$, $\vq(0)=[0.2,-0.15,-0.25,0.9354]^\top$, $\vw(0)=[0.03,0.02,-0.01]^\top$\,rad/s,
so the examples of consecutive blocks chain together: the output of one is the input of the next.

% ---------------------------------------------------------------------
\subsection{Block 1 --- How a rigid body turns}
\label{sec:b1}
% ---------------------------------------------------------------------

\begin{idea}
Two separate facts describe a turning body. \emph{Kinematics}: if it spins with rate $\vw$, its orientation changes
in a known way (this is pure geometry, no forces). \emph{Dynamics}: a torque changes the spin, and a body that spins
about two axes at once starts spinning about the third by itself (gyroscopic coupling).
\end{idea}

\begin{eqbox}{cPlant}{Equation (1): quaternion kinematics and Euler's rotational dynamics}
\begin{align}
\dot{\qv}&=\tfrac12\big(q_4I_3+\qv^{\times}\big)\vw, \qquad \dot q_4=-\tfrac12\,\qv^\top\vw \tag{1a}\label{eq:kin}\\
J\dot{\vw}&=-\vw^{\times}J\vw+\vu+\vd \tag{1b}\label{eq:dyn}
\end{align}
\end{eqbox}

\begin{dimtable}
\vq=[\qv^\top,q_4]^\top & attitude quaternion & \R^4 & --\\
\vw & body angular velocity & \R^3 & rad/s\\
J & inertia matrix & \R^{3\times3} & kg\,m$^2$\\
\vu & control torque (in block 2 replaced by $\sat(\vu)$) & \R^3 & N\,m\\
\vd & disturbance torque & \R^3 & N\,m\\
I_3 & $3\times3$ identity & \R^{3\times3} & --\\
\end{dimtable}

\subsubsection*{Where (1a) comes from}
Rotating first by $\vq$ and then by a tiny extra rotation $\vw\,dt$ (measured in the body frame) is the quaternion
product $\vq(t+dt)=\vq(t)\otimes[\tfrac12\vw\,dt;\ 1]$, because by \eqref{eq:quat} a small angle $\phi$ about $\bm n$ is
$[\tfrac{\phi}{2}\bm n;\ 1]$. Hence $\dot\vq=\tfrac12\,\vq\otimes[\vw;0]$. Writing out the Hamilton product
$[\bm a;a_4]\otimes[\bm b;b_4]=[a_4\bm b+b_4\bm a+\bm a\times\bm b;\ a_4b_4-\bm a^\top\bm b]$ with $\bm a=\qv$, $a_4=q_4$,
$\bm b=\vw$, $b_4=0$ gives
\[
\dot\vq=\tfrac12\begin{bmatrix}q_4\vw+\qv\times\vw\\ -\qv^\top\vw\end{bmatrix}
=\begin{bmatrix}\tfrac12(q_4I_3+\qv^\times)\vw\\ -\tfrac12\qv^\top\vw\end{bmatrix},
\]
which is (1a). It keeps $\|\vq\|=1$: $\tfrac{d}{dt}\|\vq\|^2=2\qv^\top\dot\qv+2q_4\dot q_4=\qv^\top\qv^\times\vw=0$, since
$\qv^\top(\qv\times\vw)=0$. (A numerical integrator does not keep it exactly, so every simulation re-normalises $\vq$ after
each step.)

\subsubsection*{Where (1b) comes from}
The angular momentum in the body frame is $\bm h=J\vw$. In the inertial frame Newton's law for rotation is
$\tfrac{d}{dt}\bm h_{\mathcal F_i}=$ torque. The body frame itself turns with $\vw$, and the derivative of any vector
seen from a turning frame picks up the term $\vw\times(\cdot)$:
\[
\tfrac{d}{dt}\bm h\big|_{\mathcal F_i}=J\dot\vw+\vw\times J\vw=\vu+\vd
\quad\Longrightarrow\quad J\dot\vw=-\vw^\times J\vw+\vu+\vd .
\]
The term $-\vw^\times J\vw$ is the gyroscopic coupling. It does no work: $\vw^\top(\vw^\times J\vw)=0$ by the
skew-symmetry of \eqref{eq:skew}. This is used in the proof of Theorem~1.

\begin{toybox}[one evaluation of Eq. (1) at the Part-1 initial state]
$q_4\vw=0.9354\,[0.03,0.02,-0.01]=[0.02806,\,0.01871,\,-0.00935]$ and
$\qv\times\vw=[0.0065,\,-0.0055,\,0.0085]$. Their sum times $\tfrac12$:
\[
\dot\qv=[0.01728,\ 0.00660,\ -0.00043]^\top\ \text{s}^{-1},\qquad \dot q_4=-\tfrac12\qv^\top\vw=-0.00275\ \text{s}^{-1}.
\]
Dynamics with the torque the controller will apply at $t=0$ (block~5 shows it saturates at $[-2.5,2.5,2.5]$\,N\,m) and
$\vd(0)=0$: $J\vw=[0.615,\,0.362,\,-0.095]$, $\vw\times J\vw=[0.00172,\,-0.0033,\,-0.00144]$, and
\[
\dot\vw=J^{-1}\big(-\vw\times J\vw+\sat(\vu)\big)=[-0.141,\ 0.144,\ 0.162]^\top\ \text{rad/s}^2 .
\]
Even at full torque the heavy spacecraft accelerates by only $\approx0.15$\,rad/s$^2$. Remember this number: it
decides which envelopes are feasible (complete documentation, Monte-Carlo study C).
\end{toybox}

\begin{projbox}
\code{paper\_reproduction.py}, function \code{rhs()}, evaluates (1a)--(1b) for a whole batch of simulations at once and
is integrated with RK4 at $\Delta t=1$\,ms; MATLAB \code{step6\_spacecraft\_closed\_loop.m} contains an independent
re-implementation. On the X3 we do \emph{not} integrate (1): Gazebo does, for the real quadcopter (\cref{sec:x3}).
\end{projbox}

% ---------------------------------------------------------------------
\subsection{Block 2 --- What the actuators can really do}
\label{sec:b2}
% ---------------------------------------------------------------------

\begin{idea}
Reaction wheels or rotors have a maximum torque. Whatever the controller asks for beyond it is simply not
delivered. The paper does not ignore this: it treats the missing torque as one more unknown disturbance.
\end{idea}

\begin{eqbox}{cPlant}{Equations (3)--(5): saturation and the lumped unknown torque}
\begin{align}
\sat(u_i)&=u_i+\pi_i,\qquad
\pi_i=\begin{cases}0,& |u_i|<\umax\\ \sgn(u_i)\,\umax-u_i,& |u_i|\ge\umax\end{cases}\tag{3--4}\\
J\dot\vw&=-\vw^\times J\vw+\vu+\bm T_d,\qquad \bm T_d=\vd+\bm\Pi,\quad \|\bm T_d\|\le\theta \tag{5}\label{eq:lumped}
\end{align}
\end{eqbox}

\begin{dimtable}
\umax & torque limit of each axis & \R & N\,m\\
\pi_i & torque lost on axis $i$ & \R & N\,m\\
\bm T_d & disturbance plus lost torque & \R^3 & N\,m\\
\theta & an (unknown) upper bound of $\|\bm T_d\|$ & \R & N\,m\\
\end{dimtable}

\subsubsection*{Derivation}
(3)--(4) is just $\sat(u_i)=\sgn(u_i)\min(|u_i|,\umax)$ written as ``command plus error''. Putting $\sat(\vu)=\vu+\bm\Pi$
into (1b) and collecting everything the controller does not know into $\bm T_d$ gives (5). The trick is that now the
plant looks as if it received the \emph{commanded} torque, plus one bounded unknown. The adaptive law (block~6) will
estimate a bound for it. $\theta$ is bounded because $\vd$ is bounded and because $\vu$ stays bounded (all its terms are
bounded while the state is inside the envelope).

\begin{toybox}[the torque lost at $t=0$]
Block 5 will compute $\vu(0)=[-10.537,\ 5.710,\ 10.414]^\top$\,N\,m. With $\umax=2.5$ every axis saturates:
$\sat(\vu)=[-2.5,\ 2.5,\ 2.5]$ and $\bm\Pi=\sat(\vu)-\vu=[8.037,\ -3.210,\ -7.914]$\,N\,m. At the very start the
``unknown'' torque is dominated by saturation, not by the disturbance (whose size is at most 0.187\,N\,m).
\end{toybox}

\begin{projbox}
\code{sat\_vec()} in \code{paper\_reproduction.py} clips each axis at 2.5\,N\,m; on the X3 the torque is clipped at
0.35\,N\,m and, in addition, the mixer must keep every rotor between 0 and 800\,rad/s (block 10).
\end{projbox}

% ---------------------------------------------------------------------
\subsection{Block 3 --- The finite-time envelope (FTPPF)}
\label{sec:b3}
% ---------------------------------------------------------------------

\begin{idea}
Draw a tube around zero that starts wide and narrows with time. If the state never touches the walls of the tube,
then automatically: its overshoot is smaller than the initial width, it is at least as fast as the walls close, and
its final error is smaller than the final width. The paper's new tube reaches its final width \emph{exactly} at a chosen
time $T_s$ and then stays constant, instead of only approaching it.
\end{idea}

\begin{eqbox}{cEnv}{Equations (6)--(8): performance constraint and finite-time prescribed performance function}
\begin{align}
-\rho_q(t)&<q_i(t)<\rho_q(t),\qquad -\rho_\omega(t)<\omega_i(t)<\rho_\omega(t),\qquad i=1,2,3 \tag{6, 9}\label{eq:constraint}\\
\rho(t)&=\begin{cases}\Big(\rho_0-\rho_T\big(1+\tfrac{t}{T_s}\big)\Big)\exp\!\Big(\dfrac{\kappa\,t}{t-T_s}\Big)+\rho_T,&0\le t<T_s\\[4pt]
\rho_T,& t\ge T_s\end{cases}\tag{8}\label{eq:ftppf}
\end{align}
For comparison, the classical exponential PPF is $\rho(t)=(\rho_0-\rho_\infty)e^{-\kappa t}+\rho_\infty$ (Eq.~7).
\end{eqbox}

\begin{dimtable}
\rho_0 & initial width: the largest allowed $|q_i|$ (or $|\omega_i|$) & \R & --, rad/s\\
\rho_T & final width: the guaranteed accuracy after $T_s$ & \R & --, rad/s\\
T_s & settling time of the envelope & \R & s\\
\kappa & shape: how early the envelope starts to close & \R & --\\
\end{dimtable}

\subsubsection*{Checking the three properties}
\begin{enumerate}[leftmargin=*]
  \item \emph{Start.} At $t=0$ the exponent is $0$, so $\rho(0)=(\rho_0-\rho_T)\cdot1+\rho_T=\rho_0$.
  \item \emph{Arrival exactly at $T_s$.} For $t\to T_s^-$ the denominator $t-T_s\to0^-$ while the numerator $\kappa t>0$,
        so the exponent $\to-\infty$ and $\exp(\cdot)\to0$. The bracket stays finite ($\to\rho_0-2\rho_T$), hence
        $\rho\to\rho_T$, and the second line continues it. So $\rho$ is continuous and equals $\rho_T$ from $T_s$ on.
        The exponential PPF instead only gets there as $t\to\infty$.
  \item \emph{It never increases and never dips below $\rho_T$.} If $\rho_0>2\rho_T$ (true for every envelope we use),
        the bracket $b(t)=\rho_0-\rho_T(1+t/T_s)$ is positive on $[0,T_s)$, so $\rho>\rho_T$ there. Differentiating,
        \[
        \dot\rho(t)=\Big(-\frac{\rho_T}{T_s}-b(t)\,\frac{\kappa T_s}{(t-T_s)^2}\Big)\exp\!\Big(\frac{\kappa t}{t-T_s}\Big)<0,
        \]
        because $\tfrac{d}{dt}\tfrac{\kappa t}{t-T_s}=-\tfrac{\kappa T_s}{(t-T_s)^2}$. Both terms are negative.
\end{enumerate}

\begin{toybox}[$\rho_q(2)$ with the paper's attitude envelope, $T_s=15$\,s]
$\rho_0=0.3$, $\rho_T=0.008$, $\kappa=2.5$:
\[
b(2)=0.3-0.008\Big(1+\tfrac{2}{15}\Big)=0.290933,\qquad \exp\!\Big(\tfrac{2.5\cdot2}{2-15}\Big)=e^{-0.38462}=0.680712,
\]
\[
\rho_q(2)=0.290933\times0.680712+0.008=0.206042 .
\]
The same for the rate envelope ($\rho_0=0.4$, $\rho_T=0.005$): $\rho_\omega(2)=0.394333\times0.680712+0.005=0.273428$\,rad/s.
Values along the way (from the same formula): $\rho_q=0.300,\ 0.252,\ 0.206,\ 0.091,\ 0.0245,\ 0.0099,\ 0.0080$ at
$t=0,1,2,5,8,10,12$\,s. With $T_s=10$\,s instead, $\rho_q(5)=0.0316$: a shorter $T_s$ closes the tube faster.
\end{toybox}

\begin{toybox}[finite time versus exponential, $T_s=8$\,s]
Exponential PPF with $\kappa=0.5$ at $t=8$\,s: $(0.3-0.008)e^{-4}+0.008=0.01335$, still $1.67\times\rho_T$. The FTPPF is
exactly $0.008$ at $8$\,s.
\end{toybox}

\begin{codebox}
\textbf{MATLAB:} \code{toy\_model/step1\_ftppf.m}. Part A prints the three numbers $0.290933$, $0.680712$, $0.206042$;
Part B checks $\rho(0)=\rho_0$, $\rho(T_s)=\rho_T$, monotonicity, and plots the effect of $T_s$ and $\kappa$.
Library function: \code{toy\_model/lib/ftppf.m}.
\end{codebox}

\begin{projbox}
\code{ftppf()} in \code{paper\_reproduction.py} (vectorised) and in \code{x3\_sim/controllers.py} (scalar, with a guard
that returns $\rho_T$ once the exponent is below $-50$ to avoid underflow). The paper's values are used for the
spacecraft; on the X3 the envelope is re-sized to $\rho_q:0.30\to0.02$, $\rho_\omega:2.5\to0.05$\,rad/s, $\kappa=2$,
$T_s=10$\,s (\cref{sec:b9}).
\end{projbox}

% ---------------------------------------------------------------------
\subsection{Block 4 --- Normalised errors}
\label{sec:b4}
% ---------------------------------------------------------------------

\begin{idea}
Instead of transforming the error through a complicated function (classical PPC), simply divide it by the current
half-width of the tube. The result says ``what fraction of the allowed room is used''. Inside the tube means
$|\varepsilon|<1$.
\end{idea}

\begin{eqbox}{cEnv}{Equation (11): the simple error transformation}
\begin{equation}
\veps_1=\frac{\qv(t)}{\rho_q(t)},\qquad \veps_2=\frac{\vw(t)}{\rho_\omega(t)},\qquad
\text{(9) holds}\iff |\varepsilon_{1i}(t)|<1\ \text{and}\ |\varepsilon_{2i}(t)|<1\ \ \forall i \tag{11}\label{eq:eps}
\end{equation}
\end{eqbox}

\begin{dimtable}
\veps_1=[\varepsilon_{11},\varepsilon_{12},\varepsilon_{13}]^\top & attitude error as a fraction of $\rho_q$ & \R^3 & --\\
\veps_2=[\varepsilon_{21},\varepsilon_{22},\varepsilon_{23}]^\top & rate error as a fraction of $\rho_\omega$ & \R^3 & --\\
\end{dimtable}

\subsubsection*{Why this is enough}
$\rho_q,\rho_\omega>0$ always, so dividing by them does not change signs and $|q_i|<\rho_q\iff|q_i/\rho_q|<1$. Nothing has
to be differentiated: classical PPC would need $\partial S^{-1}/\partial(\cdot)$ and $\dot\rho$ inside the controller. A
necessary condition is that the state \emph{starts} inside: $|q_i(0)|<\rho_{q0}$ and $|\omega_i(0)|<\rho_{\omega0}$
(paper Remark~7).

\begin{toybox}[Part-1 state at $t=0$ ($\rho_q=0.3$, $\rho_\omega=0.4$)]
\[
\veps_1=\frac{[0.2,\ -0.15,\ -0.25]}{0.3}=[0.6667,\ -0.5,\ -0.8333],
\]
\[
\veps_2=\frac{[0.03,\ 0.02,\ -0.01]}{0.4}=[0.075,\ 0.05,\ -0.025].
\]
All six are below 1 in magnitude: the start is admissible. Axis 3 uses 83\,\% of its attitude room.

Part-2 state with the paper's $\rho_{q0}=0.3$: $\varepsilon_{11}=0.3/0.3=1$ exactly. The constraint is violated at $t=0$
and $(1-\varepsilon_{11})^{-2}=1/0$. This is why we use $\rho_{q0}=0.36$ for Part~2 (\cref{warn:part2}); then
$\veps_1(0)=[0.833,\,-0.556,\,-0.833]$.
\end{toybox}

\begin{codebox}
\textbf{MATLAB:} \code{toy\_model/step2\_normalised\_error.m}, Part A, prints $\veps_1$, $\veps_2$ and the gains of block~5.
\end{codebox}

% ---------------------------------------------------------------------
\subsection{Block 5 --- The constrained PD-type control law}
\label{sec:b5}
% ---------------------------------------------------------------------

\begin{idea}
A spring and a damper: the PD part pulls the attitude error to zero ($-k_p\qv$, the spring) and brakes the rotation
($-k_d\vw$, the damper). On top of that, a second, \emph{variable} damper acts on the combination $\vs=\vw+c\qv$: its
strength is small in the middle of the tube and becomes huge near the wall. Finally $\that\,\sat(\vs/\varsigma)$ adds a
push of size $\that$ against the unknown torque.
\end{idea}

\begin{eqbox}{cCtrl}{Equations (12), (13) and (22): the controller}
\begin{align}
\vu&=-k_p\qv-k_d\vw-\vu_{cons} \tag{12}\label{eq:law}\\
\vu_{cons}&=\Big[\alpha(1-\veps_1)^{-2}+\beta(1-\veps_2)^{-2}\Big]\vs+\that\,\sgn(\vs),\qquad \vs=\vw+c\,\qv \tag{13}\label{eq:ucons}\\
u_{cons,i}&=\Big[\frac{\alpha}{(1-\varepsilon_{1i})^{2}}+\frac{\beta}{(1-\varepsilon_{2i})^{2}}\Big]s_i+\that\,\sat\!\Big(\frac{s_i}{\varsigma}\Big)
\quad\text{(Remark 5, chattering-free form)} \tag{22}\label{eq:ucons22}
\end{align}
\end{eqbox}

\begin{dimtable}
\vs & sliding variable, a ``target'' rate $-c\qv$ compared with $\vw$ & \R^3 & rad/s\\
g_i=\frac{\alpha}{(1-\varepsilon_{1i})^2}+\frac{\beta}{(1-\varepsilon_{2i})^2} & barrier gain of axis $i$ & \R & N\,m\,s\\
k_p,\ k_d & PD gains (40, 20) & \R & N\,m, N\,m\,s\\
\alpha,\ \beta & barrier weights (1.5, 1.2) & \R & N\,m\,s\\
c & weight of $\qv$ in $\vs$ (0.5) & \R & s$^{-1}$\\
\varsigma & boundary-layer width (0.02) & \R & rad/s\\
\end{dimtable}

\subsubsection*{How the barrier term keeps the state inside}
\begin{enumerate}[leftmargin=*]
  \item In the middle of the tube ($\veps\approx0$) the gain is $g_i\approx\alpha+\beta=2.7$, small next to $k_d=20$: the
        controller behaves like plain PD.
  \item As the state approaches the wall, $\varepsilon\to1$ and $(1-\varepsilon)^{-2}\to\infty$. The term $-g_is_i$ then
        dominates. Because $\vs=\vw+c\qv$, a large negative feedback on $s_i$ drives $\omega_i$ towards $-cq_i$, i.e.\ it makes
        the body rotate \emph{back towards zero} and caps its speed. This is Remark~2/Remark~7 of the paper.
  \item $\vs$ is the natural variable: when $\vs=0$ we have $\vw=-c\qv$, which by (1a) makes $\qv$ decay
        ($\tfrac{d}{dt}\|\qv\|^2=2\qv^\top\dot\qv=q_4\qv^\top\vw=-c\,q_4\|\qv\|^2<0$ for $q_4>0$). Pushing $\vs$ to zero
        therefore also pushes the attitude error to zero.
\end{enumerate}

\begin{warn}[A detail the paper leaves open: one-sided or two-sided barrier?]
Written literally, $(1-\varepsilon_{1i})^{-2}$ grows only as $\varepsilon_{1i}\to+1$. As $\varepsilon_{1i}\to-1$ it tends to
$1/4$. So the paper's barrier protects the upper wall of the tube much more strongly than the lower wall. A symmetric
alternative is $\alpha(1-\varepsilon_{1i}^2)^{-2}$, which grows at both walls; our first browser and Gazebo code used it.
We treat the \textbf{paper form as the reference} and report the symmetric form as a clearly labelled variant
(``symmetric barrier'') everywhere. Both are compared in the complete documentation.
\end{warn}

\begin{toybox}[the torque at $t=0$ for the Part-1 state, $\that(0)=0$]
Sliding variable: $\vs=\vw+0.5\,\qv=[0.03+0.1,\ 0.02-0.075,\ -0.01-0.125]=[0.13,\ -0.055,\ -0.135]$\,rad/s.

PD part: $-40\qv-20\vw=[-8+(-0.6),\ 6-0.4,\ 10+0.2]=[-8.6,\ 5.6,\ 10.2]$\,N\,m.

Barrier gains (paper form) with $\veps_1,\veps_2$ from block 4:
\[
g_1=\frac{1.5}{(1-0.6667)^2}+\frac{1.2}{(1-0.075)^2}=13.500+1.403=14.903,\quad g_2=1.996,\quad g_3=1.588 .
\]
$\vu_{cons}=g\odot\vs=[1.937,\ -0.110,\ -0.214]$, and
\[
\vu=[-8.6,\ 5.6,\ 10.2]-[1.937,\ -0.110,\ -0.214]=[-10.537,\ 5.710,\ 10.414]\ \text{N\,m}.
\]
Every axis exceeds 2.5\,N\,m, so $\sat(\vu)=[-2.5,\,2.5,\,2.5]$ (block 2). Note the one-sided effect: axis 3 is
\emph{closest} to its wall ($|\varepsilon_{13}|=0.833$) but gets the \emph{smallest} gain, because $\varepsilon_{13}<0$. With
the symmetric form the gains would be $[6.07,\ 3.87,\ 17.27]$ and $\vu=[-9.39,\ 5.81,\ 12.53]$.
\end{toybox}

\begin{codebox}
\textbf{MATLAB:} \code{step2\_normalised\_error.m} prints the gains $g$ of both forms and plots them over $\varepsilon$;
\code{step3\_control\_law.m} prints $\vs$, $\vu_{PD}$, $\vu_{cons}$, $\vu$ and $\sat(\vu)$ for the three controllers.
Library function: \code{toy\_model/lib/control\_law.m}.
\end{codebox}

\begin{projbox}
\code{control()} in \code{paper\_reproduction.py} implements (12) and (22) with \code{mode = "constrained"} (paper form),
\code{"symmetric"} or \code{"pd"} ($\vu_{cons}=0$). The value of $\varepsilon$ fed to the barrier is limited to $0.999$; this
guard is never active while the state is inside the envelope and only prevents a division by zero if it ever leaves.
\pyfile{117}{140}{\repro}{paper\_reproduction.py --- control(): Eq. (11), (12), (22)}
\end{projbox}

% ---------------------------------------------------------------------
\subsection{Block 6 --- The adaptive estimate}
\label{sec:b6}
% ---------------------------------------------------------------------

\begin{idea}
The controller does not know how large the disturbance plus the lost torque is. So it learns a number $\that$: it
grows while the sliding variable is not zero (``something is still pushing me away'') and is slowly pulled back by a
leakage term. $\that$ then sets the size of the push $\that\,\sat(\vs/\varsigma)$ in block 5.
\end{idea}

\begin{eqbox}{cAdapt}{Equation (14): adaptive law}
\begin{equation}
\dot{\that}=\frac{1}{\eta_1}\Big(\|\vs\|-\eta_2(t)\,\that\Big),\qquad \eta_1>0,\qquad
\int_0^\infty\eta_2(\tau)\,d\tau=\bar\eta<\infty \tag{14}\label{eq:adapt}
\end{equation}
Paper choice: $\eta_1=0.1$, $\eta_2(t)=2e^{-0.5t}$ (so $\bar\eta=2/0.5=4$), $\that(0)=0$.
\end{eqbox}

\subsubsection*{Reading the law}
Without leakage ($\eta_2=0$), $\that$ is the time integral of $\|\vs\|/\eta_1$: it keeps growing as long as $\vs\ne0$. The
leakage $-\eta_2\that$ pulls it down. The paper requires $\eta_2$ to have a \emph{finite} integral, which means the leakage
dies out with time; this is what makes the extra term in the Lyapunov inequality bounded (\cref{sec:theorem}).

\begin{toybox}[one Euler step, $\Delta t=1$\,ms, from $\that=0$]
$\|\vs(0)\|=\sqrt{0.13^2+0.055^2+0.135^2}=0.19532$. Then
$\that(1\,\text{ms})=0+0.001\times(0.19532-2\cdot0)/0.1=0.0019532$\,N\,m.
\end{toybox}

\begin{warn}[What the leakage that dies out means in practice]
With the chattering-free form (22) the sliding variable does not become exactly zero: inside the boundary layer
$\that\,\sat(s_i/\varsigma)$ is smaller than $\that$, so a small residual $\|\vs\|$ remains while the disturbance acts. Because
$\eta_2\to0$, nothing then stops $\that$ from creeping upwards slowly (``parameter drift''). We observe exactly this
in the reproduction: $\that$ is still rising at $t=20$\,s ($\approx0.43$), although the state has long converged. It is
harmless over the paper's horizon, but on the X3, which runs for 180\,s, we keep $\eta_2\ge0.2$ and $\that\le0.1$\,N\,m.
\code{step4\_adaptive\_law.m} shows the three cases side by side.
\end{warn}

\begin{projbox}
In \code{paper\_reproduction.py} the law is integrated with explicit Euler at 1\,ms after each RK4 step, with the
projection $\that\ge0$ (an estimate of a bound cannot be negative):
\pyfile{190}{192}{\repro}{paper\_reproduction.py --- adaptive law, Eq. (14)}
\end{projbox}

% ---------------------------------------------------------------------
\subsection{Theorem 1 --- Why the closed loop converges}
\label{sec:theorem}
% ---------------------------------------------------------------------

\begin{eqbox}{cCtrl}{Theorem 1 of the paper and its Lyapunov function (15)}
For the system (5) under the control (12)--(14), $\qv(t)\to0$ and $\vw(t)\to0$ as $t\to\infty$. Proof candidate:
\begin{equation}
V=(k_p+ck_d)\big(\qv^\top\qv+(1-q_4)^2\big)+\tfrac12\vw^\top J\vw+c\,\qv^\top J\vw+\tfrac{\eta_1}{2}\,\ttil^{\,2},
\qquad \ttil=\theta-\that \tag{15}\label{eq:V}
\end{equation}
\end{eqbox}

\subsubsection*{The argument in four steps}
\begin{enumerate}[leftmargin=*]
  \item \textbf{$V$ is a valid ``energy''.} Using $|c\,\qv^\top J\vw|\le c\,\lambda_{\max}(J)\|\qv\|\|\vw\|$ and
        $\tfrac12\vw^\top J\vw\ge\tfrac12\lambda_{\min}(J)\|\vw\|^2$, $V\ge\tfrac12\bm x^\top P\bm x$ with
        $\bm x=[\|\qv\|,\ \|\vw\|,\ |\ttil|]^\top$ and
        \[
        P=\begin{bmatrix}2(k_p+ck_d)&-c\lambda_{\max}(J)&0\\-c\lambda_{\max}(J)&\lambda_{\min}(J)&0\\0&0&\eta_1\end{bmatrix}.
        \]
        $P\succ0$ needs $2(k_p+ck_d)\lambda_{\min}(J)>(c\lambda_{\max}(J))^2$.
  \item \textbf{Differentiate along the motion.} Using (1a), (5), (12), the skew-symmetry of $\vw^\times$ and $\dot\ttil=-\dot\that$,
        the paper obtains (Eq.~17)
        \[
        \dot V=-ck_p\|\qv\|^2-k_d\|\vw\|^2+\vs^\top\bm T_d+\tfrac c2\big(q_4\vw+\vw^\times\qv\big)^\top J\vw
        -\vs^\top\big[\cdots\big]\vs-\|\vs\|\that-\ttil\big(\|\vs\|-\eta_2\that\big).
        \]
        The barrier term $-\vs^\top[\cdots]\vs$ is $\le0$ because its gain is positive (Remark~2) and can be dropped.
        Since $\vs^\top\bm T_d\le\|\vs\|\theta$, the terms with $\|\vs\|$ combine to $\|\vs\|(\theta-\that-\ttil)=0$.
  \item \textbf{Bound the leakage term} with $\ttil\that\le-\tfrac{2\sigma-1}{2\sigma}\ttil^2+\tfrac{\sigma}{2}\theta^2$ (Eq.~18),
        which gives $\dot V\le-\bm x^\top Q\bm x+\tfrac{\eta_2\sigma}{2}\theta^2$ with
        $Q=\diag\big(ck_p,\ k_d-c\lambda_{\max}(J),\ \eta_2\tfrac{2\sigma-1}{2\sigma}\big)$ (Eq.~19). $Q\succ0$ needs
        $c<k_d/\lambda_{\max}(J)$ and $\sigma>\tfrac12$.
  \item \textbf{Integrate and use Barbalat.} $V(t)\le V(0)-\lambda_{\min}(Q)\int_0^t\|\bm x\|^2+\tfrac{\sigma}{2}\theta^2\bar\eta$,
        so $V$ stays bounded, $\qv,\vw\in\mathcal L_2$, their derivatives are bounded, and Barbalat's lemma gives
        $\qv\to0$, $\vw\to0$.
\end{enumerate}

\begin{toybox}[the two conditions on $c$ with the paper's numbers]
$\lambda_{\min}(J)=14.27$, $\lambda_{\max}(J)=20.74$, $k_p=40$, $k_d=20$, $c=0.5$:
\[
Q:\quad k_d-c\lambda_{\max}=20-10.37=9.63>0\qquad(\text{i.e. } c<20/20.74=0.964),
\]
\[
P:\quad 2(40+10)\cdot14.27=1427>(0.5\cdot20.74)^2=107.5 .
\]
Both hold, so our unstated choice $c=0.5$ satisfies the theorem.
\end{toybox}

\begin{warn}[Three things a careful reader notices in the proof]
\begin{enumerate}[leftmargin=*]
  \item \textbf{What is proved.} Theorem 1 proves \emph{convergence to the origin}. It does not separately prove that
        $|\varepsilon_{li}(t)|<1$ is preserved for all $t$ (that the state never leaves the tube); the paper supports that
        with Remarks 2 and 7 and with simulations. We therefore phrase our own claims as ``constraint satisfaction
        observed in simulation'' and measure it.
  \item \textbf{Printing details of Eq.~(15).} The paper prints $(1-q_0^2)$ and $\tfrac{1}{\eta_1}\ttil^2$. With
        $(1-q_0)^2$ (the form of the cited reference [42]) and the weight $\tfrac{\eta_1}{2}$, differentiating gives exactly
        the terms of Eq.~(17); we write (15) in that consistent form. The conclusions do not change.
  \item \textbf{The proof uses $\sgn(\vs)$.} With the boundary layer (22) the cancellation in step 2 is exact only for
        $|s_i|\ge\varsigma$; inside the layer a residual of order $\varsigma$ remains. The practical guarantee is convergence to
        a small neighbourhood of zero, and $\that$ can drift (block~6). This matches what all our simulations show: the
        final $\|\vw\|$ is about $0.004$\,rad/s, not exactly zero.
\end{enumerate}
\end{warn}

\begin{projbox}
\code{lyapunov()} in \code{paper\_reproduction.py} evaluates the state part of (15) along the Part~2 trajectory
(figure \code{fig08\_theta\_hat\_and\_lyapunov.png}). The $\ttil$ term cannot be plotted, because the true bound $\theta$ of
$\bm T_d$ is unknown.
\end{projbox}
